\chapter{Methodology}
\label{ch:methodology}

\section{The LAS construction}
\label{sec:construction}

LAS is a lattice adaptor signature whose base is a simplified, Dilithium-style
Fiat--Shamir-with-aborts signature \cite{esgin2020post,ducas2018crystals}. All objects
live in $\Rq = \mathbb{Z}_q[X]/(X^d+1)$, the public matrix has the form
$A = [\,I_n \mid A'\,] \in \Rq^{\,n\times(n+\ell)}$, and a key pair is a short secret $r$
with its image $t = A\cdot r$. The hard relation $(Y,r')$ reuses exactly that structure,
% The hardness attribution follows the paper's own Section-3 hard-relation explanation,
% which is the M-SIS one.  The clause that stood here also credited M-LWE with "keeping the
% witness hidden": that is a security-analysis claim this project does not make, and the
% paper assigns M-LWE to the signature setting rather than to the relation.
so a statement is ``just another public key'': recovering a sufficiently short preimage of
$Y$ under $A$ is the Module-SIS problem underlying the relation's hardness
\cite{esgin2020post}. Esgin et al.\ \cite{esgin2020post}
name this statement-generation step \textsc{Gen}: sample an independent ternary vector
$r'$, compute $Y=A\,r'$, and
return $(Y,r')$. Thus \textsc{Gen} is not a further primitive but \textsc{KeyGen}
with its public key $t$ relabelled as the statement $Y$ and its secret $r$
relabelled as the witness $r'$. \emph{Key generation is shared}, so LAS is exactly the
base signature plus four adaptor operations. The implementation mirrors that split:
\texttt{relation\_gen} implements \textsc{Gen}, while \texttt{las.c} and
\texttt{las.rs} hold the adaptor layer alone.

The base signature commits to a mask $y$ as $w = A\,y$, forms a challenge
$c = H(\mathit{pk}, w, M)$, computes the response $\vecz = y + c\,r$ and rejects when
$\inftynorm{\vecz} > \gamma-\kappa$. Relative to a statement $Y$, the adaptor extension
\cite{esgin2020post} adds \textsc{PreSign}, which folds the statement into the hash,
$c = H(\mathit{pk}, w + Y, M)$, keeps the response $\hat{\vecz} = y + c\,r$, and rejects at
the \emph{tighter} bound
$\gamma-\kappa-1$; \textsc{PreVerify}, which recomputes $w' = A\hat{\vecz} - c\,t$ and checks
$c = H(\mathit{pk}, w'+Y, M)$; % NOT "$r' := y$": Algorithm 2 parses the pair that way, but this report reserves $y$ for
% the mask alone (see fig:flow's caption).  Writing the paper's local parse here made the
% sentence read as Adapt adding the masking randomness.
\textsc{Adapt}, which adds the witness $r'$,
$\vecz \leftarrow \hat{\vecz} + r'$, producing a signature the base verifier accepts; and
\textsc{Extract}, which recovers $s = \vecz - \hat{\vecz}$ (\cref{fig:lasfuncs}).

\begin{figure}[htbp]
  \centering
  \begin{tikzpicture}[
      font=\normalsize,
      box/.style={draw, rounded corners, align=center, minimum height=11mm,
                  inner sep=4pt},
      base/.style={box, text width=88mm},
      opbox/.style={box, text width=25mm, minimum height=18mm},
      note/.style={font=\normalsize\itshape, align=center},
      wire/.style={box, text width=104mm, minimum height=10mm},
      >={Stealth[]}]

    \node[base] (base) at (0mm,0mm)
      {Base signature reused unchanged\\
       \textsc{KeyGen}\qquad \textsc{Sign}\qquad \textsc{Verify}};

    \node[opbox] (ps) at (-42mm,-31mm)
      {PreSign\\locked $\hat{\sigma}$};
    \node[opbox] (pv) at (-14mm,-31mm)
      {PreVerify\\check $\hat{\sigma}$};
    \node[opbox] (ad) at (14mm,-31mm)
      {Adapt\\$\hat{\sigma}+r'$\\$\to\sigma$};
    \node[opbox] (ex) at (42mm,-31mm)
      {Extract\\$\sigma-\hat{\sigma}$\\$\to s$};

    \draw[->] (ps) -- (pv);
    \draw[->] (pv) -- (ad);
    \draw[->] (ad) -- (ex);

    \node[note] (hash) at (-28mm,-15mm) {statement $Y$\\in challenge hash};
    \draw[->, dashed] (hash) -- (ps);
    \draw[->, dashed] (hash) -- (pv);

    \node[note] (wit) at (14mm,-53mm) {secret witness $r'$};
    \draw[->, dashed] (wit) -- (ad);

    \node[wire] at (0mm,-70mm)
      {wire challenge\\
       transmit $\tilde c$; recompute $c=\textsf{SampleInBall}(\tilde c)$};
  \end{tikzpicture}
  \caption[The four functions LAS adds to the base signature]{The four functions LAS
  adds around an unchanged base signature. \textsc{PreSign} and \textsc{PreVerify} bind
  the public statement $Y$ into the challenge; \textsc{Adapt} uses the witness $r'$ to
  turn the pre-signature into an ordinary signature; and \textsc{Extract} recovers a
  witness from that ordinary signature and the matching pre-signature. The implementation
  transmits $\tilde c$ and re-derives $c$ with \textsf{SampleInBall}.}
  \label{fig:lasfuncs}
\end{figure}

\begin{equation}
  \label{eq:signbound}
  \text{accept iff } \inftynorm{\vecz} \le \gamma-\kappa,
  \qquad \gamma = \kappa\, d\,(n+\ell).
\end{equation}

The subtle design point is the bound budget. An \emph{honest} witness is ternary, so
$\inftynorm{r'}\le 1$, and because \textsc{PreSign} rejects at $\gamma-\kappa-1$ the response
after \textsc{Adapt} clears \cref{eq:signbound}; loosening it to $\gamma-\kappa$ would let
adapted signatures exceed the bound and be rejected.
\Cref{fig:flow} shows the data flow.

\begin{figure}[htbp]
  \centering
  \begin{tikzpicture}[
      node distance=7mm and 5mm,
      box/.style={draw, rounded corners, align=center, minimum height=8mm,
                  inner sep=3pt, font=\normalsize},
      ann/.style={font=\normalsize\itshape, align=center, text width=32mm},
      >={Stealth[]}]
    \node[box] (ps)  {PreSign\\$\hat{\vecz}=y+c\,r$\\$\hat{\sigma}=(\tilde c,\hat{\vecz})$};
    \node[box, right=of ps] (pv) {PreVerify};
    \node[box, right=of pv] (ad) {Adapt\\$\vecz=\hat{\vecz}+r'$};
    \node[box, right=of ad] (vf) {Verify /\\Publish $\sigma$};
    \node[box, right=of vf] (ex) {Extract\\$s=\vecz-\hat{\vecz}$};
    \draw[->] (ps) -- (pv);
    \draw[->] (pv) -- (ad);
    \draw[->] (ad) -- (vf);
    \draw[->] (vf) -- (ex);
    % statement Y enters at PreSign and PreVerify (hashed into the challenge)
    \node[ann, above=of ps] (Y) {statement $Y$\\(into the hash)};
    \draw[->, dashed] (Y) -- (ps);
    \draw[->, dashed] (Y) -| (pv);
    % witness y enters at Adapt; is revealed by Extract
    \node[ann, below=of ad] (y) {witness $r'$\\(secret)};
    \draw[->, dashed] (y) -- (ad);
    % anchored to ex's right edge so the wrapped annotation grows leftwards and cannot
    % push the picture's bounding box past \textwidth at 12pt.
    \node[ann, text width=38mm, anchor=north east] (yo) at ([yshift=-7mm]ex.south east)
      {$s$, recovered only with\\the matching pre-signature};
    \draw[->, dashed] (ex.south) -- (yo.north);
  \end{tikzpicture}
  \caption[The LAS adaptor data flow]{The LAS adaptor data flow, in the symbols of
  \cref{fig:lasfuncs}: $y$ is PreSign's mask, $r'$ the honest witness, $s$ what Extract
  recovers. The statement $Y$ is folded into the challenge
  hash during PreSign and PreVerify; the witness is added by Adapt to turn the
  pre-signature into a base signature; and publishing that signature lets anyone
  holding the pre-signature run Extract to recover it. This last step is what an atomic
  swap uses.}
  \label{fig:flow}
\end{figure}

\section{Implementation strategy: reuse, do not reinvent}
\label{sec:impl}

The guiding principle is that an exotic scheme equals a basic scheme plus extra functions,
so the lattice arithmetic should be \emph{reused}, not rewritten. The implementation is
built additively on the upstream CRYSTALS-Dilithium reference
\cite{ducas2018crystals,pqcrystalsdilithium} at a pinned commit, and the same strategy is
applied again in Rust (\cref{sec:rustport}). Neither is itself ``the reference'': both are
\emph{modified, simplified} ML-DSA-style bases carrying the LAS layer of
\textcite{esgin2020post}.

Classifying every source file by how the reference is used (\cref{tab:reuse}) yields a
clean split: the entire lattice core is reused unchanged, while the adaptor scheme,
serialization, application integration, and evaluation harnesses are project-specific
additions. No upstream function is modified; the low-level Dilithium arithmetic and hash
primitives are reused unchanged. The vendored
\texttt{secp256k1-zkp} \cite{secp256k1zkp} and LaZer \cite{lazer} are unmodified too, the
latter behind a dedicated bridge.

\paragraph{Data types and the matrix product.}
LAS is a self-contained module over the reference's degree-$d$ polynomial type. Because
$A=[\,I_n\mid A'\,]$, the product $A\,v = v_{\text{top}} + A' v_{\text{bot}}$ multiplies only
$A'$.

\paragraph{Hash, challenge, and samplers.}
The random oracle $H$ is built from SHAKE256, absorbing a canonical encoding of the public
key, then the commitment ($w$ for Sign, $w+Y$ for PreSign), then the message. The challenge
sampler is the reference's \texttt{SampleInBall} at weight $\kappa$, giving
$\lVert c\rVert_1=\kappa$ and $\inftynorm{c}=1$. Two samplers are added: a ternary one for
secrets and witnesses, and a rejection sampler uniform on $[-\gamma,\gamma]$ for the mask
(\cref{app:methoddetail}).

\paragraph{Simplifications, and the modulus.}
Following \cite{esgin2020post}, the base signature omits Dilithium's hint vector,
Power2Round and high/low-bit decomposition, hashing the full commitment $w$ rather
than its high bits. This keeps the exact identity $A\vecz - c\,t = w$ available, at the cost of larger keys and
signatures. Whether those optimisations can coexist with an adaptor layer was not assumed
but \emph{tested} (\cref{sec:res-mldsa}). The modulus is reused too: the reference NTT fixes
$q=8\,380\,417\approx 2^{23}$, FIPS~204's own modulus \cite{1274601}, rather than the
$q\approx 2^{24}$ of \cite{esgin2020post} --- correctness is unaffected since
$q>2\gamma$, and only the Module-SIS/LWE margin differs. No module depends on a
mode-specific constant, so identical code runs at every parameter set.

\section{Serialization and the on-chain interface}
\label{sec:serialize}

A deployable scheme exchanges \emph{bytes}, not in-memory structs, so a bit-packed wire
encoding, a typed codec and a verify-from-bytes entry point were implemented.
Two decisions matter. Validation is \emph{asymmetric}: keys and witnesses are checked on
decode, a signature is not, so a tampered response is caught by Verify --- the entry
point the tamper test attacks. And, following FIPS~204
\cite{1274601}, a signature carries not the challenge polynomial $c$ but the digest
$\tilde c$ from which every consumer re-derives it, under the standard's $\lambda/4$ rule
(\cref{app:serialize}).

\section{Testing methodology}
\label{sec:testing}

Four test types cover progressively weaker assumptions about the environment, stated in
full in \cref{tab:tests}. Functional tests assume an honest caller and assert the whole
adaptor cycle, including the \emph{tripwire} clause --- a pre-signature must \emph{fail}
basic verification --- which is what distinguishes an adaptor signature from a signature.
Serialization and tamper tests assume a hostile caller; the known-answer test assumes a
different machine, and its digest is what the second implementation is checked against
(\cref{sec:rustport}).

\section{An independent second implementation: the Rust port}
\label{sec:rustport}

The complete scheme was implemented a second time in Rust on the same reuse principle,
layered on a fixed version of the \texttt{fips204} crate \cite{fips204crate}, with the LAS
layer added as new modules. The wire format is byte-identical to the C encoder's by
construction (\cref{fig:repostructure}).

\begin{figure}[htbp]
  \centering
  \begin{tikzpicture}[
      font=\normalsize,
      box/.style={draw, rounded corners, align=center, minimum height=11mm,
                  inner sep=4pt},
      source/.style={box, text width=50mm},
      layer/.style={box, text width=112mm, minimum height=20mm},
      evidence/.style={box, text width=34mm, minimum height=13mm},
      note/.style={font=\normalsize\itshape, align=center},
      >={Stealth[]}]

    \node[source] (csrc) at (-28mm,0mm)
      {\textbf{C reference}\\\texttt{pq-crystals/dilithium}\\
       poly, NTT, SHAKE};
    \node[source] (rsrc) at (28mm,0mm)
      {\textbf{Rust crate}\\\texttt{integritychain/fips204}\\
       NTT, SHAKE, BitPack};

    \node[note] (reuse) at (0mm,-19mm)
      {primitive code reused unchanged};
    \draw[->] (csrc.south) -- (reuse);
    \draw[->] (rsrc.south) -- (reuse);

    \node[layer] (laslayer) at (0mm,-42mm)
      {\textbf{Project LAS layer, added in both languages}\\
       \texttt{setup}, \texttt{las\_types}, \texttt{relation}\\
       \texttt{basesig}, \texttt{las}, \texttt{serialize}};
    \draw[->] (reuse) -- (laslayer.north);

    \node[note] (format) at (0mm,-63mm)
      {same LAS objects and byte format};
    \draw[->] (laslayer.south) -- (format);

    \node[evidence] (tests) at (-40mm,-84mm)
      {correctness\\KAT};
    \node[evidence] (bench) at (0mm,-84mm)
      {benchmarks\\same boundary};
    \node[evidence] (swap) at (40mm,-84mm)
      {swap driver\\UTXO study};
    \draw[->] (format) -- (tests.north);
    \draw[->] (format) -- (bench.north);
    \draw[->] (format) -- (swap.north);
  \end{tikzpicture}
  \caption[Implementation structure and C/Rust correspondence]{Implementation structure and
  C/Rust correspondence across the two implementations. C uses the vendored
  \texttt{pq-crystals/dilithium} reference \cite{pqcrystalsdilithium}, Rust uses the
  \texttt{integritychain/fips204} crate \cite{fips204crate}, and neither modifies
  upstream primitive functions. The project LAS layer is added in both languages with
  the same module split and packed object format; tests, known-answer vectors, matched
  benchmarks, and the UTXO swap driver are produced from that layer. The per-file
  classification is \cref{tab:reuse}.}
  \label{fig:repostructure}
\end{figure}

The port serves three purposes. \emph{Cross-validation}: short of a formal proof, the
strongest practical argument for correctness is a second implementation agreeing on pinned
vectors. \emph{An independent
benchmark methodology}: Criterion.rs \cite{criterionrs} shares no timing code with the C
harness, so agreement defends the C numbers against harness artefacts. \emph{Deployment
realism}: a memory-safe language is the plausible production vehicle.

To make the timings comparable the two protocol drivers mirror each other exactly ---
same repetition scheme, fixed seed, fixed message, success-path gating and rejection
statistics. One caveat qualifies the port's independence: where \cite{esgin2020post} leaves
an engineering choice open, Rust adopts the C implementation's, concretely an early-abort
rejection check. Sharing that is what makes the two languages time the same algorithm, but
it is established by inspecting the two norm checks rather than by the automated evidence,
which constrains outputs and rejection rates, not the work profile of a rejected attempt
(\cref{app:methoddetail}).

\section{Benchmarking methodology and fair comparison}
\label{sec:benchmethod}

The evaluation separates \emph{computation} cost (wall-clock time per operation) from
\emph{communication} cost (bytes transmitted or stored).

\paragraph{Measurement boundaries.}
Undocumented boundaries are a recognised benchmarking pitfall for lattice signatures
\cite{riou2026apples}, so every timing belongs to one of the three declared boundaries of
\cref{tab:tiers}: a \emph{core tier} isolating the adaptor algebra from encoding, a
\emph{packed tier} giving what an on-chain consumer pays, and, for the classical baseline
alone, the \emph{hybrid} boundary its library exposes. Comparisons are made only within one
boundary, and every base-versus-LAS comparison shows \emph{both}.


\paragraph{Implementation of record, repeatability, and machine.}
Unless a caption says otherwise, timings are the C implementation's, with the Rust results
reported separately and never averaged with them (\cref{tab:rust}). Every timing is a mean
$\pm$ sample standard deviation over the repetition scheme its caption states. All
non-random inputs are fixed, so variability comes only from rejection sampling and machine
noise, never from input selection \cite{riou2026apples}. Toolchain, hardware and per-table
commands: \cref{app:repro}.

\paragraph{Per-operation reporting, not cumulative.}
Every operation is timed independently: they are split across parties, so a combined figure
would average over costs never paid together and hide the per-operation overhead.

\paragraph{The primary comparison: LAS against its own simplified base.}
LAS is exactly its base signature plus four operations on the same code, parameters and
primitives, so every comparison holds all three fixed and varies only the adaptor layer
(justified in \cref{sec:eval-strategy}). Each adaptor operation is paired with the base
operation it mirrors --- \textsc{Adapt} against \textsc{Verify}, since it runs a
PreVerify before adding the witness --- and \textsc{Extract}, having no base analogue, is
reported separately.

\paragraph{Rejection sampling: the expected attempt count.}
Sign-class operations restart until the response passes its norm bound, so no timing claim
is meaningful without the attempt statistics. Each coefficient is accepted with probability
$\bigl(2(\gamma-\kappa)+1\bigr)/(2\gamma+1)$ \emph{independently of the secret} --- why the
response is simulatable (\cref{app:methoddetail}) --- and an attempt succeeds only if all
$(n+\ell)\,d$ pass:
\begin{equation}
  \label{eq:rejacc}
  p_{\mathrm{acc}}
  = \left(\frac{2(\gamma-\kappa)+1}{2\gamma+1}\right)^{(n+\ell)d}
  \approx \left(1-\frac{\kappa}{\gamma}\right)^{(n+\ell)d}
  \approx e^{-\kappa (n+\ell) d/\gamma} = e^{-1}\approx 36.8\%,
\end{equation}
the last step using the design choice $\gamma=\kappa\,d\,(n+\ell)$ from
\cite{esgin2020post}, made exactly so the expected number of restarts is ``about $e<3$''.
Expected attempts are $1/p_{\mathrm{acc}}$, and PreSign's bound being tighter by one, it is
\emph{expected} to restart slightly more often before any adaptor arithmetic is counted.

\paragraph{Rejection sampling: the run-validity gate.}
Because the attempt count is geometric, per-call timings are heavy-tailed and averages
over too few calls can drift several percent on rejection luck alone
\cite{riou2026apples}. The drivers therefore treat the attempt statistics as a
\emph{validity gate}: measured mean attempts must match the closed-form expectation within
five standard errors, and a failing run aborts instead of producing evidence. This is a
consistency check on the rejection rate, not a proof of sampler correctness
(\cref{app:methoddetail,app:benchalg}). A second gate precedes every timing: the
correctness contract is asserted on the objects about to be timed, so no failure path is
ever measured.

\paragraph{Parameter sensitivity and component sizes.}
To test whether the adaptor overhead scales with the lattice dimensions, the same
comparison is repeated at the four parameter sets of \cref{tab:params} (symbols in
\cref{tab:notation}). Each key and signature is decomposed into components, so the
\emph{cause} of a size difference can be identified, not merely its total stated.

\begin{table}[htbp]
  \centering
  \caption[Parameter sets used in the evaluation]{Parameter sets used in the evaluation. The \emph{LAS-2020/845 reference}
  set is \emph{paper-derived} rather than that paper's exact set: it takes the dimensions
  of \cite{esgin2020post}, a historical reference point corresponding to no ML-DSA set,
  but runs at this build's modulus, which the paper permits --- only the \emph{size} of
  $q$ matters, so the concrete value may be chosen for a fast NTT. The \emph{Simplified
  Dilithium-II/III/V} sets reuse the dimensions of Dilithium modes~2/3/5, aligning the
  reusable ML-DSA primitives and the challenge-digest strength $\lvert\tilde c\rvert$
  with ML-DSA-44/65/87 while retaining LAS's own ternary secret distribution, rejection
  bounds, exact hint-free relation and unoptimised serialization. They are engineering
  benchmark settings, \emph{not} formal NIST/ML-DSA security-equivalence claims: a
  matched dimension is not a claim that LAS \emph{is} the corresponding ML-DSA parameter
  set or inherits its security category. All lattice sets
  share the ring degree $d=256$ and modulus $q=8\,380\,417\approx 2^{23}$ inherited from
  the reused Dilithium NTT; \emph{Simplified Dilithium-III} is the project's target set.
  The classical secp256k1 ECDSA adaptor is a functionality-matched baseline at
  $\approx$128-bit classical security. \Cref{tab:notation} defines each symbol.}
  \label{tab:params}
  \small
  \setlength{\tabcolsep}{4pt}
  \begin{tabular}{@{}L{0.30\linewidth}rrrrrrr@{}}
    \toprule
    Setting & $n$ & $\ell$ & $M$ & $\kappa$ & $\gamma$ & $d$ & $q$ \\
    \midrule
    LAS-2020/845 reference            & 4 & 4 & 8  & 60 & 122\,880 & 256 & 8\,380\,417 \\
    Simplified Dilithium-II           & 4 & 4 & 8  & 39 & 79\,872  & 256 & 8\,380\,417 \\
    Simplified Dilithium-III (target) & 6 & 5 & 11 & 49 & 137\,984 & 256 & 8\,380\,417 \\
    Simplified Dilithium-V            & 8 & 7 & 15 & 60 & 230\,400 & 256 & 8\,380\,417 \\
    \midrule
    Classical (secp256k1 ECDSA adaptor)
      & \multicolumn{7}{c}{$\approx$128-bit classical (see \cref{tab:classical})} \\
    \bottomrule
  \end{tabular}
\end{table}


\paragraph{The classical baseline: functionality-matched, not security-matched.}
The second baseline is the \texttt{secp256k1-zkp} library's \texttt{ecdsa\_adaptor}
module \cite{secp256k1zkp}, implementing Fournier's construction
\cite{fournier2019adaptor}, unmodified at a pinned commit. It is a functionality-matched
``price of post-quantum'' baseline, not a same-security one; plain ECDSA is excluded, the
fair counterpart of an adaptor signature being another adaptor signature.

\section{Application methodology: the UTXO swap study}
\label{sec:stage2method}

This section defines the application protocol and the measurement boundary used for the
Stage-2 results. The protocol follows Figure~1 of Esgin, Ersoy and Erkin
\cite{esgin2020post} over two UTXO ledgers; \cref{sec:res-app} later reports the execution
evidence, while \cref{sec:res-stage2,sec:res-txstruct} report cost.

\paragraph{The venue, and why it is modelled.}
The construction \cite{esgin2020post} states its applications over ``a UTXO-based
blockchain like Bitcoin \emph{where the signature algorithm is replaced with a
lattice-based signature scheme}''. That sentence is implemented literally: a UTXO ledger
was built in which \emph{the signature algorithm is a parameter}, so one ledger serves
every configuration and any measured difference is the cryptography, not the ledger. It is
a model, not a node; the boundary is set out in \cref{app:utxomodel}.

\paragraph{The protocol that is executed.}
\Cref{fig:swapflow} is the run book implemented by the driver. Alice, written $u_1$ in the
paper notation, holds the witness $y$ and therefore moves first; Bob, written $u_2$, only
pre-signs after the proof $\pi$ and Alice's pre-signature have both verified. That proof is
part of the method, not decoration: it proves the statement has a short witness, closing
the knowledge gap that would otherwise let a malicious first party leave an unadaptable
pre-signature behind. Operationally, a pre-signature is only a wallet-to-wallet protocol
message: no UTXO is spent while $\hat{\sigma}_i$ is being exchanged. A blockchain event
occurs only when an adapted signature is placed in the transaction's ordinary signature
field and that transaction is broadcast to its own ledger. The other ledger never verifies
a foreign transaction; atomicity comes from Bob observing $\sigma_2$ on ledger~B, extracting
the witness, and using it to finish the spend on ledger~A.

\begin{figure}[htbp]
  \centering
  % Type is \normalsize throughout: nothing inside a figure may be smaller than the
  % body text.  The geometry below was re-derived for 12pt -- widths from the longest
  % line in each box, row pitch from the tallest box -- so do not shrink the font back
  % to fit a coordinate.
  % The ledger boxes say "verifies $\sigma_i$", NEVER "validates witness": this report
  % uses "witness" for the secret $y$, which is never transmitted, so that wording read
  % as the ledger checking an object it never sees.
  \begin{tikzpicture}[
      font=\normalsize,
      party/.style={draw, rounded corners=3pt, align=center, minimum width=46mm,
                    minimum height=12mm, line width=0.65pt},
      msg/.style={draw, rounded corners=3pt, fill=orange!9, draw=orange!60!black,
                  align=center, text width=68mm, minimum height=13mm,
                  inner xsep=1mm, inner ysep=1mm, line width=0.6pt},
      gate/.style={draw, rounded corners=3pt, fill=yellow!13, draw=yellow!55!black,
                   align=center, text width=48mm, minimum height=16mm,
                   inner xsep=1mm, inner ysep=1mm, line width=0.6pt},
      chain/.style={draw, rounded corners=3pt, fill=blue!7, draw=blue!55!black,
                    align=center, text width=50mm, minimum height=16mm,
                    inner xsep=1mm, inner ysep=1mm, line width=0.6pt},
      >={Stealth[]}]

    \node[party, fill=orange!10, draw=orange!60!black] (u1) at (24mm,0mm)
      {Alice $(u_1)$\\knows $y$};
    \node[party, fill=blue!8, draw=blue!60!black] (u2) at (106mm,0mm)
      {Bob $(u_2)$\\extracts later};

    \node[msg] (m1) at (65mm,-17mm)
      {\textbf{1. off-chain:} send $Y,\pi,\hat{\sigma}_1$\\
       and unsigned $\mathit{tx}_1$ template};
    \node[gate] (check) at (106mm,-36mm)
      {abort gate\\verify $\pi$\\PreVerify $\hat{\sigma}_1$};
    \draw[->, line width=0.65pt] (u1.south) |- (m1.west);
    \draw[->, line width=0.65pt] (m1.east) -| (check.north);

    \node[msg] (m2) at (65mm,-55mm)
      {\textbf{2. off-chain:} send $\hat{\sigma}_2$\\
       and unsigned $\mathit{tx}_2$ template};
    \draw[->, line width=0.65pt] (check.south) |- (m2.east);

    \node[gate, fill=green!8, draw=green!45!black] (adapt2) at (24mm,-74mm)
      {Adapt with $y$\\attach $\sigma_2$\\broadcast $\mathit{tx}_2$};
    \draw[->, line width=0.65pt] (m2.west) -| (adapt2.north);

    \node[chain] (chain2) at (65mm,-93mm)
      {UTXO ledger B\\verifies $\sigma_2$\\settles $\mathit{tx}_2$};
    \draw[->, line width=0.65pt] (adapt2.south) |- (chain2.west);

    \node[gate, fill=green!8, draw=green!45!black] (extract) at (106mm,-112mm)
      {read $\sigma_2$\\Extract $y'$\\Adapt $\hat{\sigma}_1$};
    \draw[->, line width=0.65pt] (chain2.east) -| (extract.north);

    \node[chain, fill=orange!7, draw=orange!60!black] (chain1) at (65mm,-131mm)
      {UTXO ledger A\\verifies $\sigma_1$\\settles $\mathit{tx}_1$};
    \draw[->, line width=0.65pt] (extract.south) |- (chain1.east);
  \end{tikzpicture}
  \caption[The atomic-swap protocol used in the UTXO study]{The atomic-swap protocol used
  in the UTXO study, adapted from Figure~1 of \cite{esgin2020post}. Only the two adapted
  signatures reach ledgers, inside the ordinary signature/witness field of their own
  transactions; $Y$, $\pi$, the unsigned templates and the pre-signatures are off-chain
  protocol messages, and the witness $y$ is never transmitted. Step~1 gives Bob enough
  evidence to decide whether to continue; step~2 gives Alice the pre-signature she can
  adapt with $y$; once $\sigma_2$ is visible in the settled transaction on ledger~B, Bob
  extracts $y'$ and adapts the first leg for ledger~A.
  The signed message $m_i$ is the signature hash of the unsigned template
  $\mathit{tx}_i$, not the transaction bytes themselves. \Cref{sec:res-app} reports the
  assertions that this flow passed, and \cref{sec:res-stage2} measures its cost.}
  \label{fig:swapflow}
\end{figure}
\FloatBarrier

\paragraph{What a settled transaction contains, and what the adaptor layer changes.}
Of the adaptor-protocol objects exchanged by
the swap, only the adapted signature $\sigma$ reaches a ledger, occupying the slot in
which an ordinary signature already appears --- since SegWit
\cite{bip141}, a segregated \emph{witness} field. The statement $Y$, the proof $\pi$ and
both pre-signatures stay off-chain, while the witness $y$ is never transmitted: $u_2$
derives it from the published $\sigma$ and its own pre-signature. The construction
therefore adds \emph{no adaptor-specific on-chain field} --- what \emph{scriptless}
means: the swap condition need not appear as a separate on-chain script. The resulting
transaction cost is obtained by mapping the measured signature-scheme objects onto the
fields of \cite{bip141,bip144,bip341} --- signature and public key into the witness, the
committed key hash into the output script --- accepting that mapping only once it
reproduces published reference sizes for an ordinary spend (\cref{sec:res-txstruct}).

One notational point matters once the venue is a real ledger. The swap protocol writes the
signed object as $\mathit{tx}_i$, but a transaction is not a message: a signature commits to
its \emph{signature hash}. Three objects are therefore kept distinct --- the unsigned
template exchanged off-chain, its signature hash (the message passed to \textsc{PreSign}),
and the settled transaction --- since conflating them would misreport communication cost.

\paragraph{Three configurations, and what is controlled.}
The three configurations of \cref{tab:configs} each run the swap protocol end-to-end
across two ledgers. The \emph{role-A} proof is the proof of knowledge $\pi$ the protocol
places on the wire immediately after $Y$, claiming a witness to $Y$ exists and is short.
Only the $2\to3$ step is controlled: both LAS configurations share one expansion of the
public matrix, derive all key material from one pinned seed, and prove the
\emph{identical} relation
$\exists\, r' : A r' = t' \wedge \inftynorm{r'} \le 1$ --- the hard relation of the
statement/witness pair, \emph{not} the signer's key pair, and including the norm bound,
since proving only $A r' = t'$ would fail to close the knowledge gap. This is verified
rather than assumed: a configuration whose two halves disagree on a \emph{relation
binding} each computes independently refuses to run (\cref{app:methoddetail}).
The $1\to2$ step is not controlled --- it changes the signature scheme, the relation and
whether a role-A proof exists --- so it is a whole-stack comparison, the signature-only
question being answered by \cref{sec:res-classical}.

\begin{table}[htbp]
  \centering
  \caption[The three measured configurations]{The three configurations. Only the
  $2\to3$ step is controlled: those two share one public matrix, prove the identical
  relation, and receive byte-identical inputs from one pinned seed, so their difference
  is the proof system alone. Configuration~1 carries \emph{no} role-A proof, and this
  is a finding rather than an omission: LAS needs $\pi$ because of its knowledge gap
  --- extraction returns a witness of the extended relation, which need not be ternary,
  so without $\pi$ a malicious $u_1$ could claim and leave $u_2$ holding an unadaptable
  pre-signature --- whereas an ECDSA adaptor has no such gap, since extraction recovers
  the discrete logarithm exactly. The measured \texttt{secp256k1-zkp} construction
  accordingly carries no $\pi$, and adding one would have measured a modified baseline
  rather than the library as it stands. The DLEQ proof
  carried \emph{inside} a \texttt{secp256k1-zkp} pre-signature is a different object
  and is not counted as $\pi$: its author is the pre-signer, and its claim is
  pre-signature well-formedness, not knowledge of the role-A statement witness $r'$.
  Its cost is reported in the bytes by which a pre-signature exceeds a signature.}
  \label{tab:configs}
  \small
  \begin{tabular}{@{}lL{0.27\linewidth}L{0.30\linewidth}cc@{}}
    \toprule
    \# & Signature & Role-A proof $\pi$ & Setup & Fully PQ \\
    \midrule
    1 & \cfgOneSigName{}   & \cfgOnePiName{}   & ---         & \cfgOneFullyPQ{} \\
    2 & \cfgTwoSigName{}   & \cfgTwoPiName{}   & trusted     & \cfgTwoFullyPQ{} \\
    3 & \cfgThreeSigName{} & \cfgThreePiName{} & transparent & \cfgThreeFullyPQ{} \\
    \bottomrule
  \end{tabular}
\end{table}

\paragraph{The Groth16 circuit.}
No pairing-based proof of this lattice statement was available to reuse, so the circuit is
new, and cheaper than the statement suggests: $r' \mapsto A r'$ is \emph{linear with public
coefficients}, needing no witness--witness multiplication, the expensive ingredient in a
rank-one constraint system (\cref{app:methoddetail}).

\paragraph{Measurement.}
Because gas does not exist on this venue, the axes are execution time and communication
cost, the latter counting \emph{off-chain} protocol messages, not only what reaches a
ledger. Sizes are observed from the transmitted buffers rather than computed, and these
timings sit at the \textbf{packed tier}, so they are not comparable with core-tier figures
elsewhere. One-time costs are excluded and reported separately (\cref{app:methoddetail}).
